

We describe here a general framework for building neural networks that are equivariant to 3D rotations 
of their input. 3D-rotations are described by the Special Orthogonal of degree 3, aka SO(3), 
hence we speak of SO(3)-equivariant neural networks.
This framework is applicable to any input data within a spherical volume that has a semantically-meaningful
center, such as spherical images~\citep{visani_holographic-vae_2024}, directional acoustic signals~\citep{miotello_analysis_2025}
and finally atomic point clouds, the main subject of this thesis.

We first introduce group equivariance, discuss why and how to incorporate it in machine learning, and finally
introduce SO(3)-equivariant neural networks. The contents of this chapter are largely adapted from
previous works, including~\citep{pun_learning_2024, thomas_tensor_2018, kondor_clebsch_2018, bronstein_geometric_2021}.

\section{Background}

\textbf{Group Invariance and Equivariance.} Starting with an intuitive explanation, a \textit{function}
is said to be \textit{invariant} to a certain group of transformations (e.g. 3D rotations) if applying one 
such transformation to the function's input does not change its output. \textit{Equivariance} is then a
generalization of invariance whereby, when the input is transformed by the action of a certain group
element (or rather by a matrix representation parameterized by the group element), the output of the
function is transformed by the same group element (i.e., by a matrix representation parameterized
by the same group element, but that can be different from the input's representation). In short,
equivariant functions transform the input in the same way regardless of its coordinate frame, but
do not necessarily discard the coordinate frame information, whereas invariant functions also do
the latter. We note that both of these concepts can be extended to \textit{properties} as well, 
since the property of an object can be seen as the output of a function applied to the object.
For example, the mass of a molecule remains constant when rotating it (i.e. it is invariant to 3D rotations),
whereas its dipole moment rotates alongside it (i.e. it is equivariant to 3D rotations).
Switching now to a more formal treatment, a function between two vector spaces $f: X \rightarrow Y$
is said to be equivariant to a group of transformations $\mathfrak{G}$ iff applying
any group transformation to the input space of $f$ corresponds to applying the same transformation
to the output space (i.e., via a \textit{representation} parametrized by the same group element).
Formally: $f(\mD_{X}(\mathfrak{g})\vx) = \mD_{Y}(\mathfrak{g})f(\vx) \,, \forall \vx \in X \land
\forall \mathfrak{g} \in \mathfrak{G}$. The group acts on the input and output vector spaces with
via matrix multiplication with \textit{representations} that are specific to the space (e.g. 
$\mD_{X}$ and $\mD_{Y}$) and are parameterized by the same group element $\mathfrak{g}$.
A group may have different representations, and a special one is the one that always maps to the
identity: $\mD_{Y}(\mathfrak{g})=1, \forall \mathfrak{g} \in \mathfrak{G}$; a function on whose
output space $\mathfrak{G}$ acts with the identity representation is said to be \textit{invariant}
to $\mathfrak{G}$.

\textbf{Irreducible representations: the key to building equivariant neural networks.} How are
equivariant layers generally achieved? The key is to look at the group's \textit{irreducible
representations} (irreps). These are the group's \textit{smallest} representations, so that any
possible representation can be provably decomposed into a direct sum of irreps. Therefore, the
group's irreps can be used to describe how the group elements act on any vector space.
Consequently, data itself can be seen as being composed (via direct sum, i.e. concatenation) of
features of different "types", where having a particular type means transforming - under the
group's action - according to the group's corresponding irrep. It turns out that, if data is
decomposed in such a way, it is then generally easier to work out the conditions for constructing
equivariant functions acting on it. This is easier said than done, but it has been worked out for
SO(3), the group describing 3D rotations about a fixed point~\citep{geiger_e3nn_2022,
kondor_clebsch-gordan_2018, thomas_tensor_2018}. The irreps of SO(3) are the Wigner-D matrices:
\EQ
        D^{\ell}(R) \in \mathbb{C}^{(2\ell+1)\times(2\ell+1)} \quad \mathrm{where} \,\, \ell = 0,
        1, 2, ... \,\,\mathrm{and}\,\, R \in \mathrm{SO(3)}
        \label{eq:WignerD}
\EE
Notably, $D^0(R)$ describes how \textit{invariant} features transform under rotation, and thus it
is $1$ for every rotation.\\
Thus, to construct SO(3)-Equivariant models, two things are needed: 1) a way to encode the data
into a set of features that "rotate" according to these irreps, and 2) a way to process the
resulting data equivariantly.

\section{Holographic Encoding of the Data}
\label{sec:encoding}

To perform the encoding we leverage the Spherical Harmonics. This set of functions has the desired
property that each of them "rotates" according to the Wigner-D matrices:
\EQ
	Y_{\ell m}(\theta,\phi) \overset{\text{rotation:\,} R}{\longrightarrow}\sum_{m' = -\ell}^\ell
	D^{\ell}_{m'm}(R) Y_{\ell m'}(\theta,\phi)
	\label{eq:YLMWignerD}
\EE
where $Y_{\ell m}(\theta,\phi)$ is the spherical harmonic of degree $\ell$ and order $m$ defined as

\EQ
	Y_{\ell m}(\theta,\phi) =
	\sqrt{\frac{2n+1}{4\pi}\frac{(n-m)!}{(n+m)!}}e^{im\phi}P^m_\ell(\cos\theta)
	\label{eq:YLM}
\EE
where $\ell$ is a non-negative integer ($0\leq \ell$), and $m$ is an integer within the interval
$-\ell\leq m \leq \ell$. $P^m_\ell(\cos \theta)$ is the Legendre polynomial of degree $\ell$ and
order $m$.
The spherical harmonics also form a complete and orthonormal basis for signals on a sphere. Thus, a
signal over a sphere of radius $r$ - $\rho(\theta,\phi):~{S}^2(r) \rightarrow \mathbb{R}$ - can be
encoded into a set of Fourier coefficients $\hat{F}_{\ell m}$ via the spherical Fourier transform
(SFT):
\EQ
    \hat{F}_{\ell m}=\int_0^{2\pi}  \int_0^\pi \rho(\theta,\phi) Y_{\ell m} (\theta,\phi) \sin
    \theta \,\d\theta\,\d\phi
    \label{eq:SFT}
\EE
Not only do these coefficients efficiently represent the signal without loss of information (up to
desired maximum resolution $\ell_{\mathrm{max}}$), but they also transform under SO(3) following
the Wigner-D matrices (analogously to~\ref{eq:YLMWignerD}).

The original spherical signal $\rho(\theta,\phi)$ can then be recovered (up to resolution
$\ell_{\mathrm{max}}$) via the Inverse SFT (ISFT):
\EQ
    \tilde{\rho}(\theta,\phi) = \sum_{\ell = 0}^{\ell_{\mathrm{max}}} \sum_{m=-\ell}^{\ell}
    \hat{F}_{\ell m} Y_{\ell m} (\theta,\phi)
    \label{eq:ISFT}
\EE

To encode 3D data about a predefined center ($\rho(r, \theta, \phi)~: \mathbb{R}^3 \rightarrow
\mathbb{R}$) one can simply add a complete and orthonormal radial basis function that will capture
the radial distribution of the data around the center. Following~\citep{pun_learning_2022}, we use
the Zernike radial functions $R_{\ell}^{n}(r)$, leading to the Zernike Fourier Transform (ZFT):
\EQ
    \hat{Z}_{\ell m}^{n} = \int_0^1 \int_0^{2\pi}  \int_0^\pi \rho(r, \theta, \phi) Y_{\ell m}
    (\theta,\phi) R_{\ell}^{n}(r) r^2 \sin \theta \,\d r\,\d\theta\,\d\phi
    \label{eq:ZFT_prelim}
\EE
where $n$ represents radial resolution and is capped to a maximum value $n_{\mathrm{max}}$ in
practice.

To represent point clouds, a common choice for $\rho(\vr) \equiv \rho(r,\theta,\phi)$ is a sum of
Dirac-$\delta$ functions centered at each point:
\EQ
    \rho(\vr) = \sum_{i\in \text{points}} \omega_i \delta(\rho(\vr_{i}) - \rho(\vr))
    \label{eq:dirac_delta}
\EE
Where $\omega_i$ can be used to represent point-specific features.
This choice is powerful because the forward transform has a closed-form solution that does not
require a discretization of 3D space for numerical computation. Specifically, the ZFT of a point
cloud follows:
\EQ
    \hat{Z}_{\ell m}^{n} = \sum_{i\in \text{points}} \omega_i R_{n}^{\ell}(r_i) Y_{\ell
    m}(\theta_i, \varphi_i)
    \label{eq:ZFT}
\EE
And again, the data can be reconstructed up to resolution defined by $\ell_{\mathrm{max}}$ and
$n_{\mathrm{max}}$ using the corresponding inverse transform (IZFT):
\EQ
    \tilde{\rho}(r, \theta, \varphi) = \sum_{\ell = 0}^{L} \sum_{m = -\ell}^{\ell} \sum_{n}^{N}
    \hat{Z}_{\ell m}^{n} R_{\ell}^{n}(r) Y_{\ell m} (\theta, \varphi)
    \label{eq:IZFT}
\EE

\begin{figure}[b]
    \centering
    \includegraphics[width=1.0\textwidth]{Quals/figures/encoding.png}
    \figcaption{Visualization of the Holograhic Encoding considered in this work.}{}
    \medskip
    \small
    \justifying
    \textbf{(A):} Encoding of a spherical signal. The signal is encoded via the Spherical Fourier
    Transform into Fourier coefficients indexed by degree $\ell$ and order $m$, up to a maximum
    type (aka degree) $\ell_{\mathrm{max}}$ ($\ell_{\mathrm{max}} = 2$ in this example). The
    coefficients are most naturally represented with a pyramidal structure, as the order $m$ lies
    in the range $[-\ell, \ell]$. \textbf{(B):} Encoding of a centered point cloud, specifically an
    atomic environment. In addition to $\ell$ and $m$, coefficients are indexed by radial
    coefficient $n$ as well as channel $\alpha$, where the latter is indicative of atom-level
    features (e.g. atom type); in essence, the ZFT (Eg.~\ref{eq:ZFT}) is performed once per
    channel. The radial and channel indices are conceptually combined in the \textit{multiplicity}
    of a given type $\ell$.
    \label{fig:encoding}
\end{figure}

As mentioned above, we can use $\omega_i$ to introduce atom-level features (such as atom type,
partial charge, etc.). If $\omega_i$ is a \textit{vector} of multiple features (e.g. a one-hot
encoding of atom type) indexed by its own index $\alpha$, the resulting Zernike coefficients will
also have an associated index $\alpha$. This is equivalent to "splitting" the point cloud into
multiple \textit{channels} where each point has scalar feature value $\omega_i$, and concatenating
features of the same type $\ell$, as shown in Figure~\ref{fig:encoding}B. For the purposes of
constructing equivariant operations, however, the indices $\alpha$ and $n$ are combined into a
single index $c$. The number of indices $c$ for a particular type $\ell$ is called the
\textit{multiplicity} of $\ell$ ($C_{\ell}$), following the \texttt{e3nn}
nomenclature~\citep{geiger_e3nn_2022}; the multiplicity $C_{\ell}$ is often the same for all $\ell$
and just denoted by $C$.

In this document, we refer to a collection of Zernike coefficients representing a point cloud as a
\textit{hologram}, due to the fact that Zernike projections in spherical Fourier space can be
understood as a superposition of spherical holograms of an input point cloud.\\
Instead, we refer to a generic collection of features of type $\ell$ (which can be a hologram or
the output of one of the equivariant layers described in the next section) as a \textit{steerable
tensor}, following the nomenclature of ref.~\citep{brandstetter_geometric_2022}. "Tensor" indicates
the multi-dimensional nature of the representation, and "steerable" indicates that it can be
transformed ("steered") - under the action of the group of interest, SO(3) in our case - via matrix
multiplication.

\textbf{On the Choice of Zernike.} The use of other radial bases is possible within this framework,
as long as they are complete. Orthonormality is also desirable as it ensures that each basis
encodes different information, resulting in a more efficient encoding of the coefficients.
ref.~\cite{boyd_comparing_2011} demonstrates that encoding with Zernike polynomials result in a
faster convergence compared to the radial basis functions localized at different radii, as well as
most other orthogonal harmonic bases, with the exception of Logan-Shepp. ``Faster convergence"
indicates that fewer frequencies are required to encode the same information.
ref.~\cite{uhrin_through_2021} also uses Zernike to construct invariant descriptors of atomic
environments. However, this choice is not unique and other special functions such as Bessel
functions can be used to encode the data~\cite{musaelian_learning_2022}.



\section{SO(3)-Equivariant Layers}
\label{sec:layers}

Once the data is encoded in the irreps basis, one can leverage a suite of pre-existing rules to
build learnable functions (layers) without breaking rotational equivariance, and transform the
input tensor into a new tensor composed of features within the same range of possible types. The
most important consideration is that coefficients sharing the same $\ell$ and $n$, but having
different $m$ are \textit{coupled}, meaning that they cannot be transformed independently of each
other in arbitrary fashion or else equivariance gets broken. However, invariant properties (such as
$\ell = 0$ coefficients, and magnitudes $\sum_{m} |\hat{Z}_{\ell m}^{n}|^2$) can be transformed
arbitrarily.\\
Arguably the most important equivariant operation is the Clebsch-Gordan (CG) Tensor
Product~\citep{tung_group_1985}. It is the only known operation capable of non-linearly mixing
information between features of different type $\ell$ and with different momentum index $m$. The CG
tensor product combines two features of degrees $\ell_1$ and $\ell_2$ to produce another feature of
degree $|\ell_2-\ell_1|\leq\ell_3\leq |\ell_1+\ell_2|$. Let $\vh_{\ell} \in \mathbb{R}^{2\ell + 1}$
be a generic degree $\ell$ tensor, with individual components $\hat h_{\ell m}$ for $-\ell\leq m
\leq \ell $. The CG tensor product is given by,
\EQ
\begin{aligned}
    \hat  h_{\ell_3 m_3} &= (\vh_{\ell_1} \otimes_{cg} \vh_{\ell_2})_{\ell_3 m_3} \\
    &= \sum_{m_1 = -\ell_1}^{\ell_1} \sum_{m_2 = -\ell_2}^{\ell_2} C_{(\ell_1 m_1)(\ell_2 m_2)}^{(\ell_3 m_3)} \hat h_{\ell_1 m_1} \hat h_{\ell_2 m_2}
    \label{eq:cg_tensor_product}
\end{aligned}
\EE
where $C_{(\ell_1 m_1)(\ell_2 m_2)}^{(\ell_3 m_3)}$ are the Clebsch-Gordan coefficients.

These rules and operations can then be used to build neural network layers that respect
SO(3)-Equivariance.

\textbf{SO(3)-Equivariant Linearity.} The equivariant linear layer consists of a set of linear
projections each acting on the set of features sharing the same type $\ell$. It practically equates
to $\ell_{\text{max}} + 1$ standard linear layers with no bias except for the $\ell=0$ case, and
with the consideration that all $2\ell+1$ moments of the same feature are processed together.
Formally, let $\vh_{\ell} \in \mathbb{R}^{C \times (2\ell + 1)}$ be a set of $C$ features of type
$\ell$. Then, we learn weight matrix $\mW_{\ell} \in \mathbb{R}^{C \times K}$ that linearly maps
$h_{\ell}$ to $\overline{h}_{\ell} \in \mathbb{R}^{K \times (2\ell + 1)}$; if $\ell=0$,
$\mathbf{b}_{\ell} \in\mathbb{R}^{K}$ is also learned:
\EQ
    \overline{\vh}_{\ell} = \mW_{\ell}^{T} \vh_{\ell} \,\,(+ \,\mathbf{b}_{\ell})
    \label{eq:linearity}
\EE
This layer was introduced by~\citep{kondor_clebsch-gordan_2018}, and within it reside most of the
learnable parameters.

\textbf{Efficient Tensor Product Nonlinearity.} The CG Tensor Product is used as the primary
nonlinear activation within this framework -
as originally prescribed by~\citep{kondor_clebsch-gordan_2018} - by applying it between features
within the same tensor. Crucially, this is
also the only operation in this framework that can transfer information across features of
different types $\ell$.
~\citep{kondor_clebsch-gordan_2018} and ~\citep{pun_learning_2022} compute this tensor product
naively, by processing every possible pair of
features of types $\ell_1$ and $\ell_2$ within the tensor (dubbed Full Tensor Product).
Instead, we consider two modifications proposed by~\citep{cobb_efficient_2021} that significantly
reduce its computation time while sacrificing
little in terms of expressive power.
The first modification, which we refer to as "channel-wise computation", consists in conceptually
splitting
the steerable tensor along the channel dimension, yielding $C$ tensors of effective channel-size 1;
then, the Tensor Product is computed within
each teansor, and the output concatenated along the channel dimension. This modification reduces
the time and space complexity with respect to
the channel size $C$ from $O(C^2)$ to $O(C)$.
The second modification, which we refer to as "sparsification", consists in dropping some of the
individual Tensor Products. Specifically, if
there already exist a Tensor Product within ($\ell_1$, $\ell_2$), and within ($\ell_2$ , $\ell_3$),
then the Tensor Product within $(\ell_1, \ell_3)$
is dropped. This modification reduces the time and space complexity with respect to
$\ell_{\mathrm{max}}$ $C$ from $O(\ell_{\mathrm{max}}^2)$ to $O(\ell_{\mathrm{max}})$.
Throughout this work, either both modifications or only the "channel-wise computation" modification
are considered.
An ablation study on the efficacy of these modifications - specific to H-(V)AE - can be found in
Table~\ref{table:tp_ablation}. See ref.~\citep{cobb_efficient_2021} for further details.

\textbf{SO(3)-Equivariant Batch Norm.} Introduced by~\citep{kondor_clebsch-gordan_2018}. Normalizes
a tensor's features equivariantly by their respective batch-averaged magnitude:
\EQ
    N_{\ell}^{c} = \frac{1}{B} \sum_{b=1}^{B} \frac{1}{2\ell+1} \sum_{m=-\ell}^{\ell} (h_{\ell
    m}^{c b})^2
    \label{eq:batch_norm}
\EE
where $h$ is the tensor, indexed by $\ell$, $m$, $c$ and batch $b$.

\textbf{Signal Norm.} We devise a simple normalization strategy to help confine the activations
within the networks, which have a tendency to explode due to the tensor product being akin to a
square non-linearity. It consist of simply normalizing a tensor $h$ by its total magnitude.
Specifically, the total magnitude is defined as:
\EQ
    N_{tot} = \sum_{\ell} \frac{\sum_{c} \sum_{m = -\ell}^{\ell} (\hat h_{\ell m}^{c})^2}{2\ell+1}
    \label{eq:signal_norm}
\EE
Then, the tensor features are updated as $\overline{\hat h_{\ell m}^{c}} = {\hat h_{\ell m}^{c}}
w_{\ell}/{\sqrt{N_{tot}}}$, where $w_{\ell}$ is a degree-specific affine transformation for added
flexibility

\textbf{Layer Norm Nonlinearity.} This consists of applying a standard Layer
Norm~\citep{ba_layer_2016} to each of the \textit{norms} of type-$\ell$ features (which are
invariant), and then feeding the normalized norm into a standard nonlinear activation function.
This layer effectively combines the equivariant layer norm in
\texttt{e3nn}~\citep{geiger_e3nn_2022} with the Norm Nonlinearity originally used in Harmonic
Networks~\cite{worrall_harmonic_2017}, and then adapted to the SO(3) domain by Tensor Field
Networks~\citep{thomas_tensor_2018}. It is also used in the SE(3)-Transformer
~\citep{fuchs_se3-transformers_2020, noauthor_se3-transformers_nodate}.

\textbf{Implementation.} These SO(3)-equivariant layers, as well as the holographic encoding, were
implemented with the help of \texttt{e3nn} primitives~\citep{geiger_e3nn_2022}.


\section{Why Equivariance?}

In deep learning, the creation of a model that is invariant or equivariant to a desired group
of transformations has been traditionally achieved by training the network with group-transformed
copies of the input-output pairs, a practice known as data augmentation. However, for high-dimensional
groups, the amount of data augmentation required to learn the desired symmetry can be prohibitive
($\sim 500$ fold for SO(3)~\citep{geiger_e3nn_2022}). Therefore, encoding the desired symmetry structure
into the neural network structure itself, rather than learning it via data augmentation, can significantly
reduce training time while guaranteeing exactness.

When the target function is specifically invariant, another alternative is possible. Namely, 
the input may be represented using exclusively invariant features, which any un-constrained neural network
would process into an invariant output. For example, an SO(3)-invariant representation of 3 points may be
constituted by their relative distances and angles. This approach is widely used for representing atomic point
clouds due to its implementation and conceptual simplicity~\citep{dauparas_robust_2022, schutt_schnet_2017}.
However, its success relies upon constructing a representation that contains all possible invariant information
that is relevant for the task, avoiding ambiguities. For example, considering only relative distances in
the above example of 3 points would ignore key information (e.g., the angles) that is required to represent the
3 points unambiguously. Instead, using an equivariant neural network to learn an invariant output allows
for the effective learning of invariant features from the raw input, maximizing expressivity, and is recommended
in~\citep{bronstein_geometric_2021}.
Last, in the specific domain of predicting quantum properties of atomic systems, it has been empirically observed
that equivariant neural networks have more favorable scaling laws than invariant networks~\citep{batzner_e3-equivariant_2022}.




